\chapter{Boolean Analysis — Core Definitions}
%%%%%%%%%%%%%%%%%%%%%%%%%%%%%%%%%%%%%%%%%%%%%%%%%%%%%%%%%%%%%%%%%%%%%%%%%%%%%%%
\section{Overview}

This chapter develops the Fourier--Walsh theory of Boolean functions
$f : \{0,1\}^n \to \bbr$.  The uniform probability measure on $\{0,1\}^n$
endows the space of Boolean functions with an $L^2$ inner product
$\langle f, g\rangle = \E[fg]$.  The Walsh characters
$\chi_S(x) = \prod_{i\in S}(-1)^{x_i}$, indexed by subsets $S\subseteq[n]$,
form an orthonormal basis for this space, yielding the \textbf{Walsh--Fourier
expansion} $f = \sum_S \hat f(S)\,\chi_S$.

The main results are:
\begin{itemize}
  \item Walsh expansion (Fourier inversion).
  \item Parseval's identity: $\|f\|^2 = \sum_S \hat f(S)^2$.
  \item Influence via Fourier: $\mathrm{Inf}_i[f] = \sum_{S\ni i}\hat f(S)^2$.
  \item Total influence: $I[f] = \sum_S |S|\,\hat f(S)^2$.
  \item Noise operator in the Fourier domain: $\widehat{T_\rho f}(S) = \rho^{|S|}\hat f(S)$.
\end{itemize}

\textbf{Reference:} Ryan O'Donnell, \emph{Analysis of Boolean Functions},
Cambridge University Press, 2014.

%%%%%%%%%%%%%%%%%%%%%%%%%%%%%%%%%%%%%%%%%%%%%%%%%%%%%%%%%%%%%%%%%%%%%%%%%%%%%%%
\section{The Boolean hypercube and functions}

\begin{definition}[Boolean hypercube]
\lean{BooleanAnalysis.BoolCube}
\leanok
The \emph{Boolean hypercube} of dimension $n$ is $\{0,1\}^n$,
represented as the function type $\mathrm{Fin}\,n \to \mathrm{Bool}$.
\end{definition}

\begin{definition}[Boolean function]
\lean{BooleanAnalysis.BooleanFunc}
\leanok
\uses{BooleanAnalysis.BoolCube}
A \emph{Boolean function} of arity $n$ is a function $f : \{0,1\}^n \to \bbr$.
\end{definition}

%%%%%%%%%%%%%%%%%%%%%%%%%%%%%%%%%%%%%%%%%%%%%%%%%%%%%%%%%%%%%%%%%%%%%%%%%%%%%%%
\section{Uniform measure and expectation}

\begin{definition}[Uniform weight]
\lean{BooleanAnalysis.uniformWeight}
\leanok
The uniform probability measure on $\{0,1\}^n$ assigns weight
$2^{-n}$ to each point.
\end{definition}

\begin{definition}[Expectation]
\lean{BooleanAnalysis.expect}
\leanok
\uses{BooleanAnalysis.uniformWeight, BooleanAnalysis.BooleanFunc}
The \emph{expectation} of $f$ under the uniform measure:
\[
  \E[f] \;=\; 2^{-n}\sum_{x\in\{0,1\}^n} f(x).
\]
\end{definition}

\begin{definition}[$L^2$ inner product]
\lean{BooleanAnalysis.innerProduct}
\leanok
\uses{BooleanAnalysis.expect}
The \emph{$L^2$ inner product} of two Boolean functions with respect to the
uniform measure:
\[
  \langle f, g\rangle \;=\; \E[f\cdot g]
  \;=\; 2^{-n}\sum_{x\in\{0,1\}^n} f(x)\,g(x).
\]
\end{definition}

\begin{definition}[$L^2$ norm]
\lean{BooleanAnalysis.l2Norm}
\leanok
\uses{BooleanAnalysis.innerProduct}
The \emph{$L^2$ norm} of $f$: $\|f\| = \sqrt{\langle f,f\rangle}$.
\end{definition}

\begin{lemma}[Inner product is symmetric]
\lean{BooleanAnalysis.innerProduct_comm}
\leanok
\uses{BooleanAnalysis.innerProduct}
$\langle f, g\rangle = \langle g, f\rangle$.
\end{lemma}

\begin{lemma}[Inner product is nonneg on the diagonal]
\lean{BooleanAnalysis.innerProduct_self_nonneg}
\leanok
\uses{BooleanAnalysis.innerProduct}
$\langle f, f\rangle \ge 0$.
\end{lemma}

\begin{lemma}[Inner product is linear in the first argument]
\lean{BooleanAnalysis.innerProduct_add_left}
\leanok
\uses{BooleanAnalysis.innerProduct}
$\langle f + g, h\rangle = \langle f, h\rangle + \langle g, h\rangle$.
\end{lemma}

%%%%%%%%%%%%%%%%%%%%%%%%%%%%%%%%%%%%%%%%%%%%%%%%%%%%%%%%%%%%%%%%%%%%%%%%%%%%%%%
\section{Walsh--Fourier characters}

\begin{definition}[Sign encoding]
\lean{BooleanAnalysis.boolToSign}
\leanok
The \emph{sign encoding} maps $\mathrm{Bool}$ to $\{-1,1\}\subset\bbr$:
\[
  \sigma(b) \;=\; \begin{cases} 1 & b = \texttt{false},\\
                                -1 & b = \texttt{true}. \end{cases}
\]
\end{definition}

\begin{lemma}[Basic sign identities]
\lean{BooleanAnalysis.boolToSign_false}
\lean{BooleanAnalysis.boolToSign_true}
\lean{BooleanAnalysis.boolToSign_sq}
\lean{BooleanAnalysis.boolToSign_mul_self}
\lean{BooleanAnalysis.boolToSign_not}
\leanok
\uses{BooleanAnalysis.boolToSign}
$\sigma(\texttt{false}) = 1$, $\sigma(\texttt{true}) = -1$,
$\sigma(b)^2 = 1$, and $\sigma(\lnot b) = -\sigma(b)$.
\end{lemma}

\begin{definition}[Walsh character]
\lean{BooleanAnalysis.chiS}
\leanok
\uses{BooleanAnalysis.boolToSign}
The \emph{Walsh--Fourier character} associated to $S\subseteq[n]$ is
\[
  \chi_S(x) \;=\; \prod_{i\in S}(-1)^{x_i}
  \;=\; \prod_{i\in S}\sigma(x_i).
\]
The family $\{\chi_S\}_{S\subseteq[n]}$ forms an orthonormal basis for
$L^2(\{0,1\}^n,\mathrm{uniform})$.
\end{definition}

\begin{lemma}[Character of the empty set]
\lean{BooleanAnalysis.chiS_empty}
\leanok
\uses{BooleanAnalysis.chiS}
$\chi_\varnothing \equiv 1$.
\end{lemma}

\begin{lemma}[Character of a singleton]
\lean{BooleanAnalysis.chiS_singleton}
\leanok
\uses{BooleanAnalysis.chiS, BooleanAnalysis.boolToSign}
$\chi_{\{i\}}(x) = \sigma(x_i)$.
\end{lemma}

\begin{lemma}[Characters take values in $\{-1,1\}$]
\lean{BooleanAnalysis.chiS_sq_eq_one}
\leanok
\uses{BooleanAnalysis.chiS, BooleanAnalysis.boolToSign_sq}
$\chi_S(x)^2 = 1$ for all $S$ and $x$.
\end{lemma}

\begin{lemma}[Characters are multiplicative on disjoint sets]
\lean{BooleanAnalysis.chiS_mul_of_disjoint}
\leanok
\uses{BooleanAnalysis.chiS}
If $S$ and $T$ are disjoint then $\chi_{S\cup T}(x) = \chi_S(x)\cdot\chi_T(x)$.
\end{lemma}

\begin{lemma}[Product of two characters]
\lean{BooleanAnalysis.chiS_mul_chiS}
\leanok
\uses{BooleanAnalysis.chiS}
$\chi_S(x)\cdot\chi_T(x) = \chi_{S\mathbin{\triangle} T}(x)$,
where $\triangle$ denotes symmetric difference.
\end{lemma}

\begin{lemma}[Character under global bit flip]
\lean{BooleanAnalysis.chiS_neg}
\leanok
\uses{BooleanAnalysis.chiS, BooleanAnalysis.boolToSign_not}
Flipping all bits multiplies the character by $(-1)^{|S|}$:
\[
  \chi_S(\lnot x) \;=\; (-1)^{|S|}\,\chi_S(x).
\]
\end{lemma}

\begin{lemma}[Count of Walsh characters]
\lean{BooleanAnalysis.card_walsh_characters}
\leanok
\uses{BooleanAnalysis.chiS}
There are exactly $2^n$ Walsh characters, matching the dimension of
$L^2(\{0,1\}^n)$.
\end{lemma}

%%%%%%%%%%%%%%%%%%%%%%%%%%%%%%%%%%%%%%%%%%%%%%%%%%%%%%%%%%%%%%%%%%%%%%%%%%%%%%%
\section{Fourier coefficients}

\begin{definition}[Fourier coefficient]
\lean{BooleanAnalysis.fourierCoeff}
\leanok
\uses{BooleanAnalysis.innerProduct, BooleanAnalysis.chiS}
The \emph{Fourier--Walsh coefficient} of $f$ at frequency $S$ is
\[
  \hat f(S) \;=\; \langle f, \chi_S\rangle
  \;=\; 2^{-n}\sum_{x\in\{0,1\}^n} f(x)\,\chi_S(x).
\]
\end{definition}

\begin{lemma}[Fourier coefficient at $\varnothing$ equals expectation]
\lean{BooleanAnalysis.fourierCoeff_empty}
\leanok
\uses{BooleanAnalysis.fourierCoeff, BooleanAnalysis.expect}
$\hat f(\varnothing) = \E[f]$.
\end{lemma}

%%%%%%%%%%%%%%%%%%%%%%%%%%%%%%%%%%%%%%%%%%%%%%%%%%%%%%%%%%%%%%%%%%%%%%%%%%%%%%%
\section{Walsh expansion}

\begin{theorem}[Walsh expansion]
\lean{BooleanAnalysis.walsh_expansion}
\leanok
\uses{BooleanAnalysis.fourierCoeff, BooleanAnalysis.chiS}
Every Boolean function $f:\{0,1\}^n\to\bbr$ has the unique representation
\[
  f(x) \;=\; \sum_{S\subseteq[n]} \hat f(S)\,\chi_S(x).
\]
This is the Fourier inversion formula for the uniform measure on $\{0,1\}^n$.
\end{theorem}

%%%%%%%%%%%%%%%%%%%%%%%%%%%%%%%%%%%%%%%%%%%%%%%%%%%%%%%%%%%%%%%%%%%%%%%%%%%%%%%
\section{Orthonormality and Parseval's identity}

\begin{theorem}[Orthonormality of Walsh characters]
\lean{BooleanAnalysis.fourier_coeff_chi}
\leanok
\uses{BooleanAnalysis.innerProduct, BooleanAnalysis.chiS, BooleanAnalysis.chiS_mul_chiS}
The Walsh characters form an orthonormal system:
\[
  \langle \chi_S, \chi_T\rangle \;=\; [S = T].
\]
\end{theorem}

\begin{lemma}[Self inner product of a character is 1]
\lean{BooleanAnalysis.innerProduct_chi_self}
\leanok
\uses{BooleanAnalysis.fourier_coeff_chi}
$\langle \chi_S, \chi_S\rangle = 1$.
\end{lemma}

\begin{theorem}[Parseval's identity]
\lean{BooleanAnalysis.parseval}
\leanok
\uses{BooleanAnalysis.innerProduct, BooleanAnalysis.fourierCoeff,
      BooleanAnalysis.fourier_coeff_chi, BooleanAnalysis.walsh_expansion}
The squared $L^2$ norm of $f$ equals the sum of squared Fourier coefficients:
\[
  \|f\|^2 \;=\; \langle f, f\rangle \;=\; \sum_{S\subseteq[n]} \hat f(S)^2.
\]
\end{theorem}

%%%%%%%%%%%%%%%%%%%%%%%%%%%%%%%%%%%%%%%%%%%%%%%%%%%%%%%%%%%%%%%%%%%%%%%%%%%%%%%
\section{Coordinate influence}

\begin{definition}[Bit flip]
\lean{BooleanAnalysis.flipBit}
\leanok
\uses{BooleanAnalysis.BoolCube}
The point $x^i\in\{0,1\}^n$ obtained from $x$ by flipping the $i$-th coordinate.
\end{definition}

\begin{definition}[Influence]
\lean{BooleanAnalysis.influence}
\leanok
\uses{BooleanAnalysis.expect, BooleanAnalysis.flipBit}
The \emph{influence} of coordinate $i$ on $f$ measures how often flipping
bit $i$ changes the output:
\[
  \mathrm{Inf}_i[f]
  \;=\; \E\!\left[\frac{(f(x)-f(x^i))^2}{4}\right].
\]
For $\pm 1$-valued $f$ this equals $\Pr[f(x)\ne f(x^i)]$.
\end{definition}

\begin{definition}[Total influence]
\lean{BooleanAnalysis.totalInfluence}
\leanok
\uses{BooleanAnalysis.influence}
$I[f] = \sum_{i=1}^n \mathrm{Inf}_i[f]$.
\end{definition}

\begin{lemma}[Influence of a Walsh character]
\lean{BooleanAnalysis.influence_chi}
\leanok
\uses{BooleanAnalysis.influence, BooleanAnalysis.chiS}
$\mathrm{Inf}_i[\chi_S] = [i\in S]$.
\end{lemma}

\begin{theorem}[Influence via Fourier coefficients]
\lean{BooleanAnalysis.influence_eq_sum_fourier}
\leanok
\uses{BooleanAnalysis.influence, BooleanAnalysis.fourierCoeff,
      BooleanAnalysis.parseval, BooleanAnalysis.walsh_expansion,
      BooleanAnalysis.fourier_coeff_chi}
\[
  \mathrm{Inf}_i[f] \;=\; \sum_{\substack{S\subseteq[n]\\ i\in S}} \hat f(S)^2.
\]
\end{theorem}

\begin{theorem}[Total influence via Fourier coefficients]
\lean{BooleanAnalysis.totalInfluence_eq_sum_sq_deg}
\leanok
\uses{BooleanAnalysis.totalInfluence, BooleanAnalysis.influence_eq_sum_fourier}
\[
  I[f] \;=\; \sum_{S\subseteq[n]} |S|\,\hat f(S)^2.
\]
\end{theorem}

\begin{lemma}[Influence is at most 1]
\lean{BooleanAnalysis.influence_le_one}
\leanok
\uses{BooleanAnalysis.influence}
For any $\pm 1$-valued function and any coordinate $i$,
$\mathrm{Inf}_i[f] \le 1$.
\end{lemma}

\begin{lemma}[Average influence lower bound]
\lean{BooleanAnalysis.max_influence_lower_bound}
\leanok
\uses{BooleanAnalysis.totalInfluence, BooleanAnalysis.influence}
There exists a coordinate $i$ with $\mathrm{Inf}_i[f] \ge I[f]/n$.
\end{lemma}

%%%%%%%%%%%%%%%%%%%%%%%%%%%%%%%%%%%%%%%%%%%%%%%%%%%%%%%%%%%%%%%%%%%%%%%%%%%%%%%
\section{Noise operator}

\begin{definition}[Noise operator]
\lean{BooleanAnalysis.noiseOp}
\leanok
\uses{BooleanAnalysis.fourierCoeff, BooleanAnalysis.chiS}
The \emph{noise operator} $T_\rho$ with parameter $\rho\in[-1,1]$, defined via
the Fourier domain by
\[
  T_\rho f(x) \;=\; \sum_{S\subseteq[n]} \rho^{|S|}\,\hat f(S)\,\chi_S(x).
\]
Probabilistically, $T_\rho f(x) = \E_y[f(y)]$ where each bit of $y$ agrees
with $x_i$ with probability $\tfrac{1+\rho}{2}$ and is flipped with probability
$\tfrac{1-\rho}{2}$, independently.
\end{definition}

\begin{theorem}[Noise operator in the Fourier domain]
\lean{BooleanAnalysis.noiseOp_fourier}
\leanok
\uses{BooleanAnalysis.noiseOp, BooleanAnalysis.fourierCoeff,
      BooleanAnalysis.fourier_coeff_chi}
\[
  \widehat{T_\rho f}(S) \;=\; \rho^{|S|}\,\hat f(S).
\]
\end{theorem}

\begin{theorem}[Stability formula]
\lean{BooleanAnalysis.stability_formula}
\leanok
\uses{BooleanAnalysis.noiseOp, BooleanAnalysis.innerProduct,
      BooleanAnalysis.fourierCoeff, BooleanAnalysis.fourier_coeff_chi,
      BooleanAnalysis.walsh_expansion}
\[
  \langle f, T_\rho f\rangle \;=\; \sum_{S\subseteq[n]} \rho^{|S|}\,\hat f(S)^2.
\]
\end{theorem}

%%%%%%%%%%%%%%%%%%%%%%%%%%%%%%%%%%%%%%%%%%%%%%%%%%%%%%%%%%%%%%%%%%%%%%%%%%%%%%%
\section{Notable Boolean functions}

\begin{definition}[Dictator]
\lean{BooleanAnalysis.dictator}
\leanok
\uses{BooleanAnalysis.boolToSign}
The \emph{dictator function} for coordinate $i$ outputs the sign of the $i$-th
bit: $\mathrm{dict}_i(x) = (-1)^{x_i}$.  Its unique nonzero Fourier coefficient
is $\widehat{\mathrm{dict}_i}(\{i\}) = 1$.
\end{definition}

\begin{lemma}[Dictator equals singleton character]
\lean{BooleanAnalysis.dictator_eq_chi}
\leanok
\uses{BooleanAnalysis.dictator, BooleanAnalysis.chiS}
$\mathrm{dict}_i = \chi_{\{i\}}$.
\end{lemma}

\begin{definition}[Parity]
\lean{BooleanAnalysis.parity}
\leanok
\uses{BooleanAnalysis.boolToSign}
The \emph{parity function} $f(x) = (-1)^{x_1+\cdots+x_n}$ satisfies
$\hat f([n]) = 1$ and $\hat f(S) = 0$ for $S\ne[n]$.
\end{definition}

\begin{lemma}[Parity equals full character]
\lean{BooleanAnalysis.parity_eq_chi_univ}
\leanok
\uses{BooleanAnalysis.parity, BooleanAnalysis.chiS}
$\mathrm{parity}_n = \chi_{[n]}$.
\end{lemma}

\begin{definition}[Majority]
\lean{BooleanAnalysis.majority}
\leanok
\uses{BooleanAnalysis.BooleanFunc}
The \emph{majority function} on $\{0,1\}^{2k+1}$ outputs $1$ if more than half
the inputs are $\texttt{false}$, and $-1$ otherwise.
\end{definition}

%%%%%%%%%%%%%%%%%%%%%%%%%%%%%%%%%%%%%%%%%%%%%%%%%%%%%%%%%%%%%%%%%%%%%%%%%%%%%%%
\section{Sensitivity}

\begin{definition}[Sensitivity at a point]
\lean{BooleanAnalysis.sensitivity}
\leanok
\uses{BooleanAnalysis.flipBit}
The \emph{sensitivity} of $f$ at $x$ is the number of coordinates $i$ such
that flipping bit $i$ changes $f(x)$.
\end{definition}

\begin{definition}[Maximum sensitivity]
\lean{BooleanAnalysis.maxSensitivity}
\leanok
\uses{BooleanAnalysis.sensitivity}
The maximum sensitivity of $f$ over all inputs.
\end{definition}

%%%%%%%%%%%%%%%%%%%%%%%%%%%%%%%%%%%%%%%%%%%%%%%%%%%%%%%%%%%%%%%%%%%%%%%%%%%%%%%
\section{Properties for Arrow's theorem}

\begin{definition}[Odd function]
\lean{BooleanAnalysis.isOddFunc}
\leanok
\uses{BooleanAnalysis.BooleanFunc}
A Boolean function $f$ is \emph{odd} if $f(\lnot x) = -f(x)$ for all $x$.
This models antisymmetry of pairwise preferences in social choice.
\end{definition}

\begin{definition}[$\pm 1$-valued function]
\lean{BooleanAnalysis.isPmOne}
\leanok
\uses{BooleanAnalysis.BooleanFunc}
$f$ is \emph{$\pm 1$-valued} if $f(x)\in\{-1,1\}$ for all $x$.
\end{definition}

\begin{lemma}[Odd functions have zero even-level Fourier coefficients]
\lean{BooleanAnalysis.fourierCoeff_odd_even}
\leanok
\uses{BooleanAnalysis.isOddFunc, BooleanAnalysis.fourierCoeff,
      BooleanAnalysis.chiS_neg}
If $f$ is odd and $|S|$ is even, then $\hat f(S) = 0$.
\end{lemma}

\begin{lemma}[Self inner product of a $\pm 1$ function]
\lean{BooleanAnalysis.innerProduct_self_pm_one}
\leanok
\uses{BooleanAnalysis.isPmOne, BooleanAnalysis.innerProduct}
If $f$ is $\pm 1$-valued then $\langle f, f\rangle = 1$.
\end{lemma}

\begin{lemma}[Parseval for $\pm 1$-valued functions]
\lean{BooleanAnalysis.parseval_pm_one}
\leanok
\uses{BooleanAnalysis.isPmOne, BooleanAnalysis.parseval,
      BooleanAnalysis.innerProduct_self_pm_one}
If $f$ is $\pm 1$-valued then $\sum_{S\subseteq[n]}\hat f(S)^2 = 1$.
\end{lemma}

\begin{definition}[Fourier weight at level $k$]
\lean{BooleanAnalysis.weightLevel}
\leanok
\uses{BooleanAnalysis.fourierCoeff}
The \emph{Fourier weight at level $k$} is
\[
  W_k[f] \;=\; \sum_{\substack{S\subseteq[n]\\ |S|=k}} \hat f(S)^2.
\]
\end{definition}
